% ==================== 致谢 ====================
\section*{致　谢}

感谢施洪龙老师在本课程中的讲授。作业一里关于晶体结构、空间群和对称元素的训练，让我在看 PDF 卡片和指标化衍射峰时有了可以对照的模型；作业二把课堂上的布拉格方程、PDF 卡片和 Origin 绘图对到了真实的谱线上。

感谢实验指导教师和研究生师兄在仪器操作、样品制备和安全规范方面的指导。D-MAX 3 衍射仪的开机、升压、校读和关机步骤，以及“样品要压平、量要足、门要关严”这些提醒，都是数据能测出来的前提。感谢同组同学在制样、装样和记录上的配合，也感谢全班同学共享 8 组 TiO$_2$ 样品的 XRD 数据，使本文可以在同一套条件下比较不同热处理温度的谱线。

最后感谢家人和朋友在课程学习期间的支持。文中的错误和不足由本人负责。
